\documentclass[11pt]{article}
\usepackage{amsmath,amssymb,amsthm}
\usepackage{hyperref}
\usepackage[margin=1in]{geometry}

\theoremstyle{plain}
\newtheorem{theorem}{Theorem}
\newtheorem{conjecture}[theorem]{Conjecture}
\newtheorem{proposition}[theorem]{Proposition}
\newtheorem{lemma}[theorem]{Lemma}
\newtheorem{corollary}[theorem]{Corollary}
\theoremstyle{definition}
\newtheorem{definition}[theorem]{Definition}
\newtheorem{example}[theorem]{Example}
\newtheorem{remark}[theorem]{Remark}
\newtheorem{problem}[theorem]{Problem}

\newcommand{\N}{\mathbb{N}}
\newcommand{\Z}{\mathbb{Z}}
\newcommand{\Q}{\mathbb{Q}}
\newcommand{\GK}{\operatorname{GK}}
\newcommand{\AG}{\operatorname{AG}}
\newcommand{\CPP}{\operatorname{CPP}}
\newcommand{\qpoch}[2]{(#1;#2)_\infty}
\newcommand{\qpochn}[3]{(#1;#2)_{#3}}
\newcommand{\qbin}[2]{\genfrac{[}{]}{0pt}{}{#1}{#2}}
\newcommand{\jtheta}[2]{\theta(#1;#2)}

\title{The Positivity Conjecture for $Q_{n,c}(q)$}
\author{Reference summary for research on cylindric partitions}
\date{\today}

\begin{document}
\maketitle

\begin{abstract}
This document collects the statement, context, and motivation for the
positivity conjecture on the polynomials $Q_{n,c}(q)$ arising from the
two-variable generating function of rank-3 cylindric partitions.
The conjecture first appears as Conjecture~4.2 of Corteel--Dousse--Uncu
\cite{CDU} and independently as Conjecture~2.7 of Warnaar \cite{W}.
We give the relevant definitions from the theory of cylindric partitions,
state the conjecture precisely, and illustrate it with examples.
\end{abstract}

\tableofcontents

%% ============================================================
\section{Motivation: from Rogers--Ramanujan to $\mathrm{A}_2$}
\label{sec:motivation}
%% ============================================================

\subsection{The Rogers--Ramanujan identities}

The Rogers--Ramanujan identities \cite{RR}
\begin{align}
  \sum_{n \geq 0} \frac{q^{n^2}}{\qpochn{q}{q}{n}}
  &= \prod_{n \geq 0} \frac{1}{(1-q^{5n+1})(1-q^{5n+4})},
  \label{eq:RR1} \\
  \sum_{n \geq 0} \frac{q^{n^2+n}}{\qpochn{q}{q}{n}}
  &= \prod_{n \geq 0} \frac{1}{(1-q^{5n+2})(1-q^{5n+3})},
  \label{eq:RR2}
\end{align}
are among the most celebrated $q$-series identities.
They have interpretations in the theory of partitions, the representation
theory of affine Lie algebras, conformal field theory, and statistical
mechanics.

\subsection{The Andrews--Gordon identities}

Gordon \cite{Gordon} generalised the combinatorial form of the
Rogers--Ramanujan identities to arbitrary odd moduli $2k+1$.
Andrews \cite{Andrews} found the analytic counterpart: for
$r \geq 2$ and $1 \leq s \leq r$,
\begin{equation}\label{eq:AG}
  \sum_{\substack{n_1 \geq n_2 \geq \cdots \geq n_{r-1} \geq 0}}
  \frac{q^{n_1^2 + \cdots + n_{r-1}^2 + n_s + \cdots + n_{r-1}}}
       {\qpochn{q}{q}{n_1} \prod_{i=1}^{r-2} \qbin{n_i}{n_{i+1}}}
  = \frac{\qpoch{q^{2r+1}}{q^{2r+1}}}{\qpoch{q}{q}}
    \,\jtheta{q^s}{q^{2r+1}}.
\end{equation}
The cases $r = 2$, $s \in \{1, 2\}$ recover \eqref{eq:RR1}--\eqref{eq:RR2}.
These identities correspond to principally specialised characters of the
affine Lie algebra $\mathrm{A}_1^{(1)}$.

\subsection{The $\mathrm{A}_2$ problem}

Finding higher-rank generalisations of the Andrews--Gordon identities
for $\mathrm{A}_{r-1}^{(1)}$ with $r \geq 3$ has been a long-standing open problem.
Andrews, Schilling and Warnaar \cite{ASW} discovered several
$\mathrm{A}_2$ Rogers--Ramanujan identities (moduli 4 and 7) using an
$\mathrm{A}_2$ Bailey lemma, but the resulting $q$-series were expressed as
characters multiplied by a factor of $(q;q)_\infty$ that obscured positivity.

The breakthrough insight, due to Foda--Welsh \cite{FW} and
Corteel \cite{Corteel}, was that \textbf{cylindric partitions} provide
a natural combinatorial framework for these identities.  The key
observation is that the generating function for cylindric partitions of
rank~3, when evaluated at $z = 1$, yields the $\mathrm{A}_2$
Andrews--Gordon characters (up to a universal factor).

This raises a fundamental question: can one find \emph{manifestly positive}
multisum expressions for all $\mathrm{A}_2$ Andrews--Gordon identities?
Corteel, Dousse and Uncu \cite{CDU} made significant progress on this
for level $d = 5$, and the general positivity conjecture emerged from
their work and that of Warnaar \cite{W}.


%% ============================================================
\section{Cylindric partitions}
\label{sec:cpp}
%% ============================================================

\subsection{Definition}

\begin{definition}[Cylindric partition, rank 3]
\label{def:cpp}
Let $c = (c_0, c_1, c_2)$ be a composition of $d := c_0 + c_1 + c_2$.
A \emph{cylindric plane partition of profile~$c$} is a triple of partitions
$\nu = (\nu^{(1)}, \nu^{(2)}, \nu^{(3)})$ satisfying the interlacing conditions:
\begin{align}
  \nu^{(1)}_j &\geq \nu^{(2)}_{j+c_1} && \text{for all } j \geq 1, \label{eq:interlace1} \\
  \nu^{(2)}_j &\geq \nu^{(3)}_{j+c_2} && \text{for all } j \geq 1, \label{eq:interlace2} \\
  \nu^{(3)}_j &\geq \nu^{(1)}_{j+c_0} && \text{for all } j \geq 1. \label{eq:interlace3}
\end{align}
Here $\nu^{(i)}_j = 0$ whenever $j$ exceeds the length of $\nu^{(i)}$.
The \emph{size} and \emph{max} statistics are
\[
  |\nu| := |\nu^{(1)}| + |\nu^{(2)}| + |\nu^{(3)}|,
  \qquad
  \max(\nu) := \max\bigl(\nu^{(1)}_1,\, \nu^{(2)}_1,\, \nu^{(3)}_1\bigr).
\]
\end{definition}

\begin{remark}
This is a specialisation of the general definition of Gessel--Krattenthaler
\cite{GK} to rank $r = 3$.  The profile $c$ encodes the ``shifts'' in the
interlacing conditions; equivalently, the cylindric partition lives on a
cylinder of circumference~$3$ with profile determining the shape of the
fundamental domain.
\end{remark}

\begin{example}[Profile $(2, 1)$]
\label{ex:profile-21}
For $c = (2, 1, 0)$ (so $d = 3$), a cylindric partition is a triple
$(\nu^{(1)}, \nu^{(2)}, \nu^{(3)})$ where:
\begin{itemize}
  \item $\nu^{(1)}_j \geq \nu^{(2)}_{j+1}$ (shift by $c_1 = 1$),
  \item $\nu^{(2)}_j \geq \nu^{(3)}_j$ (shift by $c_2 = 0$),
  \item $\nu^{(3)}_j \geq \nu^{(1)}_{j+2}$ (shift by $c_0 = 2$).
\end{itemize}
\end{example}

\subsection{The generating functions}

\begin{definition}[Two-variable generating function]
\label{def:GK}
The bivariate generating function for cylindric partitions with profile~$c$ is
\begin{equation}
  \GK_c(z,q) := \sum_{\nu \in \CPP(c)} z^{\max(\nu)}\, q^{|\nu|},
  \label{eq:GK-def}
\end{equation}
where the sum runs over all cylindric partitions of profile~$c$.
\end{definition}

\begin{definition}[Size generating function]
\label{def:GK-size}
The (univariate) size generating function is obtained by setting $z = 1$:
\[
  \GK_c(q) := \GK_c(1,q) = \sum_{\nu \in \CPP(c)} q^{|\nu|}.
\]
\end{definition}

\subsection{Borodin's product formula}

\begin{theorem}[Borodin \cite{Borodin}]
\label{thm:borodin}
The size generating function factors as an infinite product.  For rank~3
with profile $c = (c_0, c_1, c_2)$ and $d = c_0 + c_1 + c_2$,
\[
  \GK_c(q) = \frac{1}{\qpoch{q}{q}} \cdot
  \frac{1}{\displaystyle\prod_{\square \in \mu} \qpoch{q^{h(\square)}}{q^d}
  \cdot \prod_{\square \in \mu^c} \qpoch{q^{d-h(\square)}}{q^d}},
\]
where $\mu = \mu(c)$ is the partition associated to the profile, $h(\square)$
is the hook length, and $\mu^c$ is the complement of $\mu$ in the
$3 \times s(1,3)$ rectangle, with $s(i,j) = c_i + c_{i+1} + \cdots + c_{j-1}$.
\end{theorem}

\subsection{Connection to $\mathrm{A}_2$ characters}

The crucial link to representation theory is the identity (Foda--Welsh \cite{FW}):
\begin{equation}
  \AG_{\lambda,3}(q) = \qpoch{q}{q} \cdot \GK_c(q),
  \label{eq:AG-GK}
\end{equation}
where $\AG_{\lambda,3}(q)$ is the $\mathrm{A}_2$ Andrews--Gordon series associated
to the $\mathrm{A}_2^{(1)}$ weight
$\lambda = (d - \mu_1 + \mu_3)\Lambda_0 + (\mu_1 - \mu_2)\Lambda_1 + (\mu_2 - \mu_3)\Lambda_2$,
and the profile entries are
$c_0 = d - \mu_1 + \mu_3$, $c_i = \mu_i - \mu_{i+1}$ for $1 \leq i \leq 2$.
This means that every identity about $\GK_c(q)$ translates to a statement
about specialised affine Lie algebra characters.


%% ============================================================
\section{The $E$--$H$ factorisation and the polynomials $Q_{n,c}(q)$}
\label{sec:EH}
%% ============================================================

\subsection{Factoring out the universal series}

The generating function $\GK_c(z,q)$ contains a ``universal'' factor
independent of the profile. Defining the corrected generating function
\begin{equation}
  G_c(z,q) := \qpoch{zq}{q} \cdot \GK_c(z,q),
  \label{eq:Gc-def}
\end{equation}
the factor $H(z) = 1/\qpoch{zq}{q} = \sum_{n \geq 0} q^n z^n/\qpochn{q}{q}{n}$
(representing ordinary partitions with largest part $\leq n$, or
equivalently with at most~$n$ parts; see Section~\ref{sec:why-z-max}
below) is stripped off.
Note that $\GK_c(z,q) = H(z) \cdot G_c(z,q)$ is the $E$--$H$
factorisation of the cylindric partition generating function.

\subsection{Why $z$ tracks the maximum entry}
\label{sec:why-z-max}

The variable~$z$ in $\GK_c(z,q) = \sum z^{\max(\mathfrak{c})}\, q^{|\mathfrak{c}|}$
tracks the maximum entry $\max(\mathfrak{c}) = \max_k \mu^k_1$.  This is
natural because the universal factor $H(z)$ admits two equivalent readings,
related by the conjugation involution $\lambda \mapsto \lambda'$ on
partitions:
\begin{equation}\label{eq:Hz-two-readings}
  H(z) = \frac{1}{\qpoch{zq}{q}}
  = \sum_\lambda z^{\lambda_1}\, q^{|\lambda|}
  = \sum_\lambda z^{\ell(\lambda)}\, q^{|\lambda|},
\end{equation}
where the first sum tracks the \emph{largest part} $\lambda_1$ and the
second tracks the \emph{number of parts} $\ell(\lambda)$.  These are the
same generating function because conjugation preserves size and swaps
$\lambda_1 \leftrightarrow \ell(\lambda)$.

\begin{itemize}
  \item \textbf{Height reading} ($z^{\lambda_1}$): the $z$-degree
    measures the height of the tallest column.  This is the statistic
    that appears in the definition of $\GK_c(z,q)$.
  \item \textbf{Depth reading} ($z^{\ell(\lambda)}$): the $z$-degree
    measures the number of parts, which in the growth diagram equals the
    number of peeling layers needed to decompose the CPP.
\end{itemize}

The $E$--$H$ factorisation then expresses $z^{\max(\mathfrak{c})}$ as a
convolution of the base partition's contribution (height or depth of~$\gamma$)
and the purely cylindric contribution from $G_c(z,q)$.

\subsection{The corrected coefficients}

Write the $z$-expansion of the corrected generating function as
\begin{equation}
  G_c(z,q) = \sum_{n \geq 0} g_{n,c}(q)\, z^n,
  \label{eq:gn-def}
\end{equation}
where $g_{n,c}(q) = [z^n]\,G_c(z,q)$.  By the Euler expansion
$\qpoch{zq}{q} = \sum_{j \geq 0} (-1)^j q^{j(j+1)/2} z^j / \qpochn{q}{q}{j}$,
we have the alternating sum expression:
\begin{equation}
  g_{n,c}(q)
  = \sum_{j=0}^{n} \frac{(-1)^j\, q^{j(j+1)/2}}{\qpochn{q}{q}{j}}
    \cdot [z^{n-j}]\,\GK_c(z,q).
  \label{eq:gn-alt-sum}
\end{equation}

\subsection{The polynomial $Q_{n,c}(q)$}

\begin{definition}[$Q_{n,c}(q)$]
\label{def:Q}
For a nonneg integer $n$ and profile $c = (c_0, c_1, c_2)$ with
$d = c_0 + c_1 + c_2$ and $\ell = \gcd(d, 3)$, define
\begin{equation}
  Q_{n,c}(q) := \qpochn{q^\ell}{q^\ell}{n}
  \cdot [z^n]\!\Bigl(\qpoch{zq}{q} \cdot \GK_c(z,q)\Bigr)
  = \qpochn{q^\ell}{q^\ell}{n} \cdot g_{n,c}(q).
  \label{eq:Q-def}
\end{equation}
When $d \not\equiv 0 \pmod{3}$, we have $\ell = 1$ and
$Q_{n,c}(q) = \qpochn{q}{q}{n} \cdot g_{n,c}(q)$.
\end{definition}

\begin{remark}
The multiplication by $\qpochn{q^\ell}{q^\ell}{n}$ clears the denominators
appearing in \eqref{eq:gn-alt-sum}, converting the formal power series
$g_{n,c}(q)$ into a polynomial.
That $Q_{n,c}(q)$ is indeed a polynomial (and not merely a formal series)
was proved by Welsh \cite{Welsh}.
\end{remark}

\subsection{The $q$-binomial decomposition}

The definition yields a useful expansion in terms of $q$-binomial coefficients.
Write $f_{m,c}(q) := [z^m]\,\GK_c(z,q)$ for the $z$-coefficients of the
\emph{original} (uncorrected) generating function, and define
\begin{equation}
  h_{m,c}(q) := \qpochn{q}{q}{m} \cdot f_{m,c}(q).
  \label{eq:hm-def}
\end{equation}
Note that $h_{m,c}(q) \neq Q_{m,c}(q)$: the former uses the coefficients
$f_{m,c}$ of $\GK_c(z,q)$, while the latter uses the coefficients
$g_{m,c}$ of the corrected $G_c(z,q) = \qpoch{zq}{q} \cdot \GK_c(z,q)$.

Then:
\begin{equation}
  Q_{n,c}(q) = \sum_{j=0}^{n} (-1)^j\, q^{j(j+1)/2}\,
  \qbin{n}{j} \cdot h_{n-j,c}(q).
  \label{eq:Q-qbin}
\end{equation}
This follows from expanding $\qpoch{zq}{q}$ by the Euler identity
$\qpoch{zq}{q} = \sum_j (-1)^j q^{j(j+1)/2} z^j / \qpochn{q}{q}{j}$,
extracting $[z^n]$, and multiplying through by $\qpochn{q}{q}{n}$,
using $\qpochn{q}{q}{n}/\qpochn{q}{q}{j} = \qbin{n}{j} \cdot \qpochn{q}{q}{n-j}$.

Computational evidence shows that $h_{m,c}(q)$ is itself a polynomial with
nonneg coefficients for all tested profiles.  This is a \emph{stronger}
observation than positivity of $Q_{n,c}(q)$, since $Q_{n,c}(q)$ is built
from the $h_{m,c}(q)$ via an alternating-sign sum.


%% ============================================================
\section{The conjecture}
\label{sec:conjecture}
%% ============================================================

We now state the positivity conjecture.  The formulation below combines
the statements of Corteel--Dousse--Uncu \cite[Conjecture~4.2]{CDU}
and Warnaar \cite[Conjecture~2.7]{W}.

\begin{conjecture}[Corteel--Dousse--Uncu; Warnaar]
\label{conj:main}
Let $c = (c_0, c_1, c_2)$ and $d := c_0 + c_1 + c_2$ with
$d \not\equiv 0 \pmod{3}$.  Then:
\begin{enumerate}
  \item[(i)] $Q_{n,c}(q)$ is a polynomial in $q$ with
    \textbf{nonneg coefficients}.
  \item[(ii)] The evaluation at $q = 1$ is
    \begin{equation}
      Q_{n,c}(1)
      = \left(\frac{(d+1)(d+2)}{6} - 1\right)^{\!n}.
      \label{eq:Q-at-1}
    \end{equation}
\end{enumerate}
\end{conjecture}

\begin{remark}[Status]
The polynomiality of $Q_{n,c}(q)$ and the $q = 1$ evaluation \eqref{eq:Q-at-1}
have been \textbf{proved} by Welsh \cite{Welsh}.
The positivity (i) \textbf{remains open} in general.
\end{remark}

\begin{remark}[The $d \equiv 0 \pmod{3}$ case]
When $3 \mid d$ the appropriate normalisation uses $\ell = 3$, and the
polynomial $P_{n,c}(q) := \qpochn{q^3}{q^3}{n} \cdot g_{n,c}(q)$ is
conjectured to have nonneg coefficients as well
(Corteel--Dousse--Uncu \cite[Conjecture~4.2]{CDU}).
\end{remark}

\begin{remark}[Equivalent formulation via $g_{n,c}$]
Writing $g_{n,c}(q) = Q_{n,c}(q)/\qpochn{q}{q}{n}$ (when $\ell = 1$),
the conjecture says that the generating function of cylindric partitions
with max entry exactly~$n$, after correction by $\qpoch{zq}{q}$, has the
form ``nonneg polynomial divided by $(q)_n$''.  In other words, the
alternating-sign cancellation in \eqref{eq:gn-alt-sum} resolves perfectly
into positivity after clearing denominators.
\end{remark}


%% ============================================================
\section{Examples}
\label{sec:examples}
%% ============================================================

\subsection{Level $d = 2$: the Rogers--Ramanujan identities}

For $d = 2$ there are two profiles (up to the shift symmetry
$\GK_{(c_0,c_1,c_2)}(z,q) = \GK_{(c_1,c_2,c_0)}(z,q)$):

\paragraph{Profile $c = (2,0,0)$.}
The interlacing conditions force $\nu^{(2)} = \nu^{(3)} = \emptyset$ and
$\nu^{(1)}$ is any partition with parts $\leq n$ (when $\max = n$).
Here
\[
  Q_{n,(2,0,0)}(q) = q^{n^2},
\]
a single monomial.  Combining with the $E$--$H$ factorisation and Borodin's
formula reproduces the \textbf{first Rogers--Ramanujan identity} \eqref{eq:RR1}:
\[
  \GK_{(2,0,0)}(1,q) = \frac{1}{\qpoch{q}{q}} \cdot
  \frac{1}{\qpoch{q^2}{q^5} \cdot \qpoch{q^3}{q^5}},
  \qquad
  \AG_{\Lambda_0+\Lambda_1;3}(q) = \sum_{n \geq 0} \frac{q^{n^2}}{\qpochn{q}{q}{n}}.
\]

\paragraph{Profile $c = (1,1,0)$.}
Here
\[
  Q_{n,(1,1,0)}(q) = q^{n(n+1)},
\]
again a single monomial, and the identity becomes the
\textbf{second Rogers--Ramanujan identity} \eqref{eq:RR2}:
\[
  \AG_{2\Lambda_0;3}(q) = \sum_{n \geq 0} \frac{q^{n^2+n}}{\qpochn{q}{q}{n}}.
\]

In both cases the conjecture holds trivially since a monomial has
nonneg coefficients.


\subsection{Level $d = 4$: the Andrews--Schilling--Warnaar identities}

For $d = 4$ we have $\frac{1}{6}(5)(6) - 1 = 4$, so $Q_{n,c}(1) = 4^n$.

\paragraph{Profile $c = (2,2,0)$.}
The first few polynomials are:
\begin{align*}
  Q_{0,(2,2,0)}(q) &= 1, \\
  Q_{1,(2,2,0)}(q) &= q^2 + q^3 + q^4 + q^5, \\
  Q_{2,(2,2,0)}(q) &= q^4 + q^5 + 3q^6 + 3q^7 + 5q^8 + 4q^9 + 5q^{10}
    + 3q^{11} + 3q^{12} + q^{13} + q^{14}.
\end{align*}
These are manifestly positive and palindromic, with $Q_{1,(2,2,0)}(1) = 4$.

\paragraph{Profile $c = (2,1,1)$.}
The generating function $\GK_{(2,2,1)}(z,q)$ was computed by
Corteel--Dousse--Uncu \cite[Theorem~3.3]{CDU}.
Their result yields:
\[
  Q_{n,(2,2,1)}(q) = \sum_{\substack{n_2, n_3, n_4 \geq 0 \\
  m_1, \ldots, m_{k-1} \geq 0}}
  q^{n^2 + n_2^2 + n_3^2 + n_4^2 - nn_2 + n_2n_4}
  \qbin{n}{n_2} \qbin{n}{n_4} \qbin{n_2}{n_3},
\]
which gives $Q_{n,(2,2,1)}(1) = 6^n$.
This is a manifestly positive quadruple sum.


\subsection{Level $d = 5$: the Corteel--Dousse--Uncu identities}

For $d = 5$ we have $\frac{1}{6}(6)(7) - 1 = 6$, so $Q_{n,c}(1) = 6^n$.
This is the main case treated in \cite{CDU}, where manifestly positive
multisum expressions for $\GK_c(z,q)$ were found for all five profiles
(up to symmetry).


\subsection{Level $d = 7$}

For $d = 7$ we have $\frac{1}{6}(8)(9) - 1 = 11$, so $Q_{n,c}(1) = 11^n$.
There are four distinct profiles up to cyclic shift: $(3,2,2)$, $(2,3,2)$,
$(4,2,1)$, and $(5,1,1)$.  Computational verification confirms positivity
of $Q_{n,c}(q)$ for all profiles with $n \leq 3$.


%% ============================================================
\section{Why cylindric partitions?}
\label{sec:why-cpp}
%% ============================================================

\subsection{The product side}

Borodin's formula (Theorem~\ref{thm:borodin}) gives a closed-form
\emph{infinite product} for $\GK_c(1,q)$.  Via equation \eqref{eq:AG-GK},
this product is (up to the factor $1/\qpoch{q}{q}$) an $\mathrm{A}_2^{(1)}$
character.  So each profile~$c$ corresponds to a specific affine Lie algebra
module, and the generating function of cylindric partitions is its character.

\subsection{The sum side}

The conjecture posits that the \emph{sum side}---the multisum expression
for $Q_{n,c}(q)$---has nonneg coefficients.  Finding these multisum
expressions amounts to finding \textbf{manifestly positive} formulas for
$\mathrm{A}_2$ Andrews--Gordon identities.

The connection is:
\begin{equation}
  \AG_{\lambda,3}(q)
  = \qpoch{q}{q} \cdot \GK_c(1,q)
  = \sum_{n \geq 0} \frac{Q_{n,c}(q)}{\qpochn{q}{q}{n}}.
  \label{eq:AG-from-Q}
\end{equation}
So positivity of $Q_{n,c}(q)$ would give a manifestly positive
``sum = product'' identity for every $\mathrm{A}_2$ Andrews--Gordon
character (at levels not divisible by~3).

\subsection{The role of the max statistic}

The key innovation is the \emph{two-variable} generating function
$\GK_c(z,q)$, which refines the character by the max statistic.
Without the variable~$z$, one only sees the total count $\GK_c(1,q)$
and the connection to $Q_{n,c}(q)$ is lost.  The max statistic provides the
$z$-grading that decomposes the character into the pieces
$g_{n,c}(q) = [z^n]\,G_c(z,q)$, each of which carries the positivity.

This was first observed by Corteel--Welsh \cite{CW}, who showed that
the Corteel--Welsh $q$-difference equation (Theorem~2.3 of \cite{CDU}) for
$\GK_c(z,q)$ can be solved recursively to produce manifestly positive
expressions when $d$ is small.


%% ============================================================
\section{The hierarchy of generating functions}
\label{sec:hierarchy}
%% ============================================================

The generating function $\GK_c(z,q)$ that appears in the conjecture is the
simplest in a hierarchy of progressively more refined generating series
for cylindric partitions.  Each level specialises to the one below it, and
the conjecture lives at the bottom of the chain.  Understanding the richer
levels may be the key to proving positivity.

\subsection{Level 0: the univariate generating function}

The coarsest invariant tracks only the total size.  For a binary string
$\pi$ of length~$T$ encoding the profile, Borodin's identity
\cite{Borodin} gives:
\begin{equation}
  \sum_{\mathfrak{c} \in \CPP(\pi)} z^{|\mathfrak{c}|}
  = \prod_{n \geq 0} \left(
    \frac{1}{1 - z^{(n+1)T}}
    \prod_{\substack{i < j \\ \pi_i > \pi_j}} \frac{1}{1 - z^{j-i+nT}}
    \prod_{\substack{i > j \\ \pi_i > \pi_j}} \frac{1}{1 - z^{j-i+(n+1)T}}
  \right).
  \label{eq:borodin-full}
\end{equation}

\subsection{Level 1: the bivariate generating function $\GK_c(z,q)$}

This refines Level~0 by tracking the max statistic alongside the size:
\[
  \GK_c(z,q) = \sum_{\mathfrak{c} \in \CPP(c)}
  z^{\max(\mathfrak{c})}\, q^{|\mathfrak{c}|}.
\]
This is the level at which the conjecture is formulated.
The variable~$z$ grades by the largest entry, decomposing $\GK_c(1,q)$
into the pieces $g_{n,c}(q) = [z^n]\,G_c(z,q)$ from which $Q_{n,c}(q)$ is
constructed.  Setting $z = 1$ recovers Level~0.


\subsection{Level 2: the multivariate refinement}

A cylindric partition $\mathfrak{c} = (\mu^0, \mu^1, \ldots, \mu^{T-1})$
consists of $T$ individual partitions.  The \emph{multivariate generating
function} (PhD thesis \cite{Langer-thesis}, Chapter~3) tracks each slice
separately with its own variable:
\begin{equation}
  \sum_{\mathfrak{c} \in \CPP(\pi)}
  z_0^{|\mu^0|}\, z_1^{|\mu^1|} \cdots z_{T-1}^{|\mu^{T-1}|}.
  \label{eq:multivariate-gf}
\end{equation}
The \emph{refined Borodin identity} states that this equals the product
\eqref{eq:borodin-full} with the substitutions:
\begin{align}
  z^{nT} &\mapsto z_0^n z_1^n \cdots z_{T-1}^n, \notag \\
  z^{j-i+nT} &\mapsto z_0^n \cdots z_i^n\, z_{i+1}^{n+1} \cdots
    z_j^{n+1}\, z_{j+1}^n \cdots z_{T-1}^n
    && (i < j), \label{eq:refined-sub} \\
  z^{j-i+(n+1)T} &\mapsto z_0^{n+1} \cdots z_j^{n+1}\, z_{j+1}^n \cdots
    z_i^n\, z_{i+1}^{n+1} \cdots z_{T-1}^{n+1}
    && (i > j). \notag
\end{align}

Setting $z_0 = z_1 = \cdots = z_{T-1} = 1$ recovers Level~0.
Level~1 ($\GK_c(z,q)$) is obtained from Level~2 by the non-linear
operation of replacing
$z_0^{|\mu^0|} \cdots z_{T-1}^{|\mu^{T-1}|}$ by $z^{\max_k \mu^k_1} q^{|\mathfrak{c}|}$.


\subsection{Level 3: the operator trace (most general Schur form)}

The proof in the PhD thesis \cite{Langer-thesis} actually establishes an
identity in \emph{free parameters} $u_0, u_1, \ldots, u_T$ and $w$, from
which the multivariate version is a specialisation.

The left-hand side is expressed as a \emph{trace} of operators acting on the
ring of symmetric functions:
\begin{equation}
  \sum_{\mathfrak{c} \in \CPP(\pi)}
  \prod_{k=0}^{T-1} z_k^{|\mu^k|}
  = \sum_\mu \langle s_\mu \mid
    G^{\pi_0}(u_0)\, G^{\pi_1}(u_1) \cdots G^{\pi_{T-1}}(u_{T-1})\,
    D_w \mid s_\mu \rangle,
  \label{eq:trace-schur}
\end{equation}
where $G^1(u) = \Omega[Xu]$ (strip addition), $G^0(u) = \Omega^*[Xu]$
(strip removal), $D_w$ is the degree operator
$D_w\, s_\lambda = w^{|\lambda|}\, s_\lambda$, and
\begin{equation}
  w = z_0 z_1 \cdots z_{T-1}, \qquad
  u_k = \begin{cases}
    z_0 z_1 \cdots z_{k-1} & \text{if } \pi_k = 1, \\
    z_0^{-1} z_1^{-1} \cdots z_{k-1}^{-1} & \text{if } \pi_k = 0.
  \end{cases}
  \label{eq:uk-specialization}
\end{equation}

The proof proceeds in three steps:
\begin{enumerate}
  \item \textbf{Commute} all $\Omega^*$ operators past all $\Omega$ operators
    using $\Omega^*[Xu]\,\Omega[Xv] = \frac{1}{1-uv}\,\Omega[Xv]\,\Omega^*[Xu]$,
    producing scalar factors $1/(1 - u_i u_j)$ for each pair with
    $\pi_i > \pi_j$ and $i < j$.
  \item \textbf{Shift} using $D_w\,\Omega[Xu] = \Omega[Xuw]\,D_w$ and iterate,
    defining
    \begin{equation}
      M_\pi(m) = \sum_\mu \langle s_\mu \mid
      \prod_{\substack{k:\,\pi_k=0}} \Omega[Xu_k w^m]
      \prod_{\substack{k:\,\pi_k=1}} \Omega^*[Xu_k]\,
      D_w \mid s_\mu \rangle,
      \label{eq:M-def}
    \end{equation}
    which satisfies the recurrence
    \begin{equation}
      M_\pi(m) = \prod_{\substack{(i,j)\\\pi_i \neq \pi_j}}
      \frac{1}{1 - u_i u_j w^{m+1}} \;\cdot\; M_\pi(m+1).
      \label{eq:M-recurrence}
    \end{equation}
  \item \textbf{Take the limit} $m \to \infty$, where
    $M_\pi(\infty) = \prod_{n \geq 1} 1/(1-w^n)$.
\end{enumerate}

The result is an identity in the free parameters $u_0, \ldots, u_{T-1}, w$:
\begin{equation}
  \operatorname{tr}\!\bigl(
  G^{\pi_0}(u_0) \cdots G^{\pi_{T-1}}(u_{T-1})\, D_w
  \bigr)
  = \prod_{m \geq 1} \frac{1}{1 - w^m}
  \prod_{\substack{i < j \\ \pi_i > \pi_j}}
    \frac{1}{1 - u_i u_j w^{m-1}}
  \prod_{\substack{i > j \\ \pi_i > \pi_j}}
    \frac{1}{1 - u_i u_j w^m}.
  \label{eq:borodin-trace}
\end{equation}
The multivariate identity (Level~2) is recovered by the specialisation
\eqref{eq:uk-specialization}, and the univariate identity (Level~0) by
the further specialisation $z_k \mapsto z$ for all~$k$.


\subsection{Level 4: the $(r,t)$-Macdonald extension}

The richest form replaces the Schur inner product by the Macdonald inner
product and each cylindric partition by a weighted object.  For
$\mathfrak{c} = (\mu^0, \mu^1, \ldots, \mu^T)$, define the
\emph{Macdonald weight}
\begin{equation}
  W_{\mathfrak{c}}(r,t)
  := \prod_{\substack{k = 1 \\ \pi_k = 1}}^{T}
     \varphi_{\mu^k / \mu^{k-1}}(r,t)
     \prod_{\substack{k = 1 \\ \pi_k = 0}}^{T}
     \psi_{\mu^{k-1} / \mu^k}(r,t),
  \label{eq:Wc-def}
\end{equation}
where $\varphi_{\lambda/\mu}(r,t)$ and $\psi_{\lambda/\mu}(r,t)$ are the
Macdonald Pieri coefficients:
\begin{align}
  \varphi_{\lambda/\mu}(r,t)
  &= \prod_{s \in C_{\lambda/\mu}}
     \frac{1 - r^{a_\lambda(s)}\, t^{\ell_\lambda(s)+1}}
          {1 - r^{a_\lambda(s)+1}\, t^{\ell_\lambda(s)}},
  \label{eq:phi-def} \\
  \psi_{\lambda/\mu}(r,t)
  &= \prod_{s \notin C_{\lambda/\mu}}
     \frac{1 - r^{a_\lambda(s)+1}\, t^{\ell_\lambda(s)}}
          {1 - r^{a_\lambda(s)}\, t^{\ell_\lambda(s)+1}},
  \label{eq:psi-def}
\end{align}
with $a_\lambda(s)$ and $\ell_\lambda(s)$ the arm and leg lengths of the
box~$s$ in~$\lambda$, and $C_{\lambda/\mu}$ the set of columns where the
skew shape $\lambda/\mu$ has its boxes.

The left-hand side of the $(r,t)$-identity is expressed as an operator trace
over Macdonald polynomials (PhD thesis \cite{Langer-thesis}, Chapter~4):
\begin{equation}
  \sum_{\mathfrak{c} \in \CPP(\pi)}
  W_{\mathfrak{c}}(r,t)
  \prod_{k=0}^{T-1} z_k^{|\mu^k|}
  = \sum_\mu \langle Q_\mu \mid
    H^{\pi_0}(u_0)\, H^{\pi_1}(u_1) \cdots H^{\pi_{T-1}}(u_{T-1})\,
    D_w \mid P_\mu \rangle_{r,t},
  \label{eq:trace-mac}
\end{equation}
where $H^1(u) = \Omega_{r,t}[Xu]$, $H^0(u) = \Omega_{r,t}^*[Xu]$, and
$\langle\,\cdot\,,\,\cdot\,\rangle_{r,t}$ is the Macdonald inner product
with respect to which $\{P_\mu\}$ and $\{Q_\mu\}$ are dual bases.
The proof is identical in structure to the Schur case, using the
$(r,t)$-commutation relation
\[
  \Omega_{r,t}^*[Xu]\;\Omega_{r,t}[Xv]
  = \frac{\qpochn{tuv}{r}{\infty}}{\qpochn{uv}{r}{\infty}}
  \;\Omega_{r,t}[Xv]\;\Omega_{r,t}^*[Xu]
\]
in place of $\Omega^*[Xu]\,\Omega[Xv] = \frac{1}{1-uv}\,\Omega[Xv]\,\Omega^*[Xu]$.

The $(r,t)$-recurrence $N_\pi(m) \to N_\pi(m+1)$ produces the factors
$\qpochn{tu_iu_jw^{m+1}}{r}{\infty}/\qpochn{u_iu_jw^{m+1}}{r}{\infty}$,
and the same limit $N_\pi(\infty) = \prod_{n \geq 1} 1/(1-w^n)$ gives:
\begin{equation}
  \operatorname{tr}\!\bigl(
  H^{\pi_0}(u_0) \cdots H^{\pi_{T-1}}(u_{T-1})\, D_w
  \bigr)_{r,t}
  = \prod_{m \geq 1} \frac{1}{1 - w^m}
  \prod_{\substack{i < j \\ \pi_i > \pi_j}}
    \frac{\qpochn{tu_iu_jw^{m-1}}{r}{\infty}}
         {\qpochn{u_iu_jw^{m-1}}{r}{\infty}}
  \prod_{\substack{i > j \\ \pi_i > \pi_j}}
    \frac{\qpochn{tu_iu_jw^m}{r}{\infty}}
         {\qpochn{u_iu_jw^m}{r}{\infty}}.
  \label{eq:rt-borodin-trace}
\end{equation}

The weight function has a combinatorial description in terms
of \emph{peak} and \emph{valley} cubes (PhD thesis, Theorem~4.2):
\begin{equation}
  W_{\mathfrak{c}}(r,t)
  = \Omega\!\left[(r - t)\,\mathcal{D}_{\mathfrak{c}}(r,t)\right],
  \label{eq:Wc-plethystic}
\end{equation}
where $\Omega[\cdot]$ is the plethystic exponential and
\[
  \mathcal{D}_{\mathfrak{c}}(r,t)
  = \sum_{s \in \mathrm{peak}(\mathfrak{c})}
    r^{a_{\mathfrak{c}}(s)}\, t^{\ell_{\mathfrak{c}}(s)}
  - \sum_{s \in \mathrm{valley}(\mathfrak{c})}
    r^{a_{\mathfrak{c}}(s)}\, t^{\ell_{\mathfrak{c}}(s)}.
\]

\subsection{Summary: the nesting}

The five levels form a chain of specialisations:

\medskip
\begin{center}
\fbox{\parbox{0.92\textwidth}{\medskip
\textbf{Level 4} --- $(r,t)$-operator trace in free parameters $u_k, w$
\emph{(most general)}:
\[
  \operatorname{tr}\!\bigl(
  H^{\pi_0}(u_0) \cdots H^{\pi_{T-1}}(u_{T-1})\, D_w
  \bigr)_{r,t}
  = \prod_{m \geq 1} \frac{1}{1-w^m}
  \prod_{\pi_i > \pi_j}
  \frac{(tu_iu_jw^{\bullet};r)_\infty}{(u_iu_jw^{\bullet};r)_\infty}
\]

\smallskip
$\Big\downarrow$ \emph{set $r = t$ (all Macdonald weights $\to 1$)}

\smallskip
\textbf{Level 3} --- Schur operator trace in free parameters $u_k, w$:
\[
  \operatorname{tr}\!\bigl(
  G^{\pi_0}(u_0) \cdots G^{\pi_{T-1}}(u_{T-1})\, D_w
  \bigr)
  = \prod_{m \geq 1} \frac{1}{1-w^m}
  \prod_{\pi_i > \pi_j}
  \frac{1}{1-u_iu_jw^{\bullet}}
\]

\smallskip
$\Big\downarrow$ \emph{specialise $u_k, w$ via \eqref{eq:uk-specialization}}

\smallskip
\textbf{Level 2} --- Multivariate: $\;\displaystyle\sum_{\mathfrak{c}}
z_0^{|\mu^0|} z_1^{|\mu^1|} \cdots z_{T-1}^{|\mu^{T-1}|}$

\smallskip
$\Big\downarrow$ \emph{set all $z_k = 1$}
\hfill
$\Big\downarrow$ \emph{max extraction $\to z^{\max_k \mu_1^k} q^{|\mathfrak{c}|}$}

\smallskip
\textbf{Level 0} --- Size: $\;\displaystyle\sum_{\mathfrak{c}} q^{|\mathfrak{c}|}$
\hfill
\textbf{Level 1} --- Bivariate: $\;\GK_c(z,q) = \displaystyle\sum_{\mathfrak{c}} z^{\max(\mathfrak{c})} q^{|\mathfrak{c}|}$

\smallskip
\hfill $\Big\downarrow$ \emph{set $z = 1$}

\smallskip
\hfill \textbf{Level 0} --- Size: $\;\displaystyle\sum_{\mathfrak{c}} q^{|\mathfrak{c}|}$

\medskip
}}
\end{center}
\medskip

Level~1 (the bivariate $\GK_c(z,q)$) is the level at which
Conjecture~\ref{conj:main} is stated.  It sits ``beside'' Level~2
rather than strictly below it, because the max statistic is obtained
from the multivariate data by a non-linear operation (taking the maximum
of the first parts).

%% ============================================================
\section{The Macdonald extension of the conjecture}
\label{sec:mac-conjecture}
%% ============================================================

\subsection{Notation}

The Macdonald polynomials $P_\lambda(X;r,t)$ depend on two parameters
traditionally called $q$ and~$t$.  Since we have already used $q$ as the
weight-tracking variable throughout this document (as in $q^{|\mathfrak{c}|}$
and $Q_{n,c}(q)$), we adopt the convention of writing $(r,t)$ for the
Macdonald parameters to avoid ambiguity.  The specialisation $r = t$
recovers the Schur function setting.

\subsection{The $(r,t)$-extended generating function}

Define the \emph{Macdonald-weighted bivariate generating function}:
\begin{equation}
  \GK_c^{(r,t)}(z,q)
  := \sum_{\mathfrak{c} \in \CPP(c)}
  W_{\mathfrak{c}}(r,t) \cdot z^{\max(\mathfrak{c})}\, q^{|\mathfrak{c}|},
  \label{eq:GK-rt}
\end{equation}
where the Macdonald weight $W_{\mathfrak{c}}(r,t)$ is as defined in
\eqref{eq:Wc-def} (with the substitution $q \mapsto r$):
\[
  W_{\mathfrak{c}}(r,t)
  = \prod_{\substack{k:\,\pi_k = 1}}
    \varphi_{\mu^k/\mu^{k-1}}(r,t)
    \prod_{\substack{k:\,\pi_k = 0}}
    \psi_{\mu^{k-1}/\mu^k}(r,t).
\]
This specialises to the ordinary generating function when $r = t$:
\[
  \GK_c^{(r,r)}(z,q) = \GK_c(z,q).
\]

\subsection{The corrected coefficients}

As in the Schur case, define the corrected generating function and
its $z$-coefficients:
\begin{align}
  G_c^{(r,t)}(z,q)
  &:= \qpoch{zq}{q} \cdot \GK_c^{(r,t)}(z,q),
  \label{eq:Gc-rt} \\
  g_{n,c}^{(r,t)}(q)
  &:= [z^n]\, G_c^{(r,t)}(z,q).
  \label{eq:gn-rt}
\end{align}
Note that the $E$--$H$ factorisation uses $\qpoch{zq}{q}$, which depends
only on the weight variable~$q$ and \emph{not} on the Macdonald
parameters $(r,t)$.  This is because the ``universal factor''
$H(z) = 1/\qpoch{zq}{q}$ counts ordinary partitions by size,
independently of any Macdonald weighting.

\subsection{The polynomial $Q_{n,c}(q;\,r,t)$}

\begin{definition}
\label{def:Q-rt}
For a nonneg integer $n$, profile $c = (c_0, c_1, c_2)$ with
$d = c_0 + c_1 + c_2$, $d \not\equiv 0 \pmod{3}$, and Macdonald
parameters $(r,t)$, define
\begin{equation}
  Q_{n,c}(q;\,r,t)
  := \qpochn{q}{q}{n} \cdot g_{n,c}^{(r,t)}(q)
  = \qpochn{q}{q}{n} \cdot [z^n]\!\Bigl(
    \qpoch{zq}{q} \cdot \GK_c^{(r,t)}(z,q)
  \Bigr).
  \label{eq:Q-rt}
\end{equation}
\end{definition}

By construction, $Q_{n,c}(q;\,r,r) = Q_{n,c}(q)$.

\begin{remark}[Structure of the coefficients]
Since the Macdonald Pieri coefficients $\varphi_{\lambda/\mu}(r,t)$ and
$\psi_{\lambda/\mu}(r,t)$ are rational functions of $r$ and~$t$,
the quantity $Q_{n,c}(q;\,r,t)$ is \emph{a priori} a formal power series
in~$q$ with coefficients in $\Q(r,t)$.  For it to be a polynomial in $q$
with ``positive'' coefficients, we would need:
\begin{enumerate}
  \item[(a)] $Q_{n,c}(q;\,r,t)$ is a polynomial in~$q$
    (finitely many nonzero $q$-powers), and
  \item[(b)] each $q$-coefficient is a nonneg expression in $r$ and~$t$.
\end{enumerate}
The precise meaning of ``nonneg in $r,t$'' requires care: the Pieri
coefficients are ratios of terms of the form $(1 - r^a t^b)$, so the
natural positivity domain is $0 < r, t < 1$.
\end{remark}

\subsection{The conjecture}

\begin{conjecture}[Macdonald extension]
\label{conj:mac}
Let $c = (c_0, c_1, c_2)$ with $d := c_0 + c_1 + c_2$ and
$d \not\equiv 0 \pmod{3}$.  Then:
\begin{enumerate}
  \item[(i)] $Q_{n,c}(q;\,r,t)$ is a polynomial in~$q$ whose
    coefficients are rational functions of $r,t$ that are
    \textbf{nonneg} for $0 < r, t < 1$.
  \item[(ii)] Specialising $r = t$ recovers Conjecture~\ref{conj:main}:
    \[
      Q_{n,c}(q;\,r,r) = Q_{n,c}(q).
    \]
  \item[(iii)] At $q = 1$:
    \[
      Q_{n,c}(1;\,r,t) = \left(\frac{(d+1)(d+2)}{6} - 1\right)^{\!n},
    \]
    independent of $(r,t)$.
\end{enumerate}
\end{conjecture}

\begin{remark}[Why (iii) should be independent of $(r,t)$]
The evaluation at $q = 1$ counts cylindric partitions at the
``combinatorial'' level (ignoring the $q$-grading), and the Macdonald
weight satisfies $W_{\mathfrak{c}}(r,t)\big|_{q=1} = 1$ for any
cylindric partition $\mathfrak{c}$, since the Pieri coefficients
collapse: the products of $(1 - r^a t^b)$ factors in numerator and
denominator cancel pairwise when the $q$-weight is forgotten.
More precisely, summing the weight variable $q^{|\mathfrak{c}|}$
at $q = 1$ eliminates the distinction between different
$q$-refinements, so the $q = 1$ evaluation depends only on the
combinatorial count.
\end{remark}

\begin{remark}[The $(r,t)$-Borodin product at $z = 1$]
Setting $z = 1$ in $\GK_c^{(r,t)}(z,q)$ gives a $(r,t)$-deformed
character.  Via the $(r,t)$-Borodin identity
\eqref{eq:rt-borodin-trace}, this is:
\[
  \GK_c^{(r,t)}(1,q) = \prod_{n \geq 1} \frac{1}{1-q^n}
  \prod_{\substack{i < j \\ \pi_i > \pi_j}}
    \frac{\qpochn{tq^{j-i+(n-1)T}}{r}{\infty}}
         {\qpochn{q^{j-i+(n-1)T}}{r}{\infty}}
  \prod_{\substack{i > j \\ \pi_i > \pi_j}}
    \frac{\qpochn{tq^{j-i+nT}}{r}{\infty}}
         {\qpochn{q^{j-i+nT}}{r}{\infty}},
\]
where $(a;r)_\infty = \prod_{k \geq 0}(1 - ar^k)$.
At $r = t$ the $(r,t)$-factors reduce to~$1$ and we recover the
ordinary Borodin product.
\end{remark}

\begin{remark}[Why the conjecture is natural]
Each individual Macdonald weight $W_{\mathfrak{c}}(r,t)$ is a product of
positive rational functions (for $0 < r, t < 1$), so each cylindric
partition contributes a positive term to $\GK_c^{(r,t)}(z,q)$.
The alternating-sign cancellation in the $E$--$H$ factorisation
(multiplication by $\qpoch{zq}{q}$) is the same Euler expansion
$\sum (-1)^j q^{j(j+1)/2} z^j / (q)_j$ as in the Schur case.
If the cancellation resolves into positivity for $r = t$ (the original
conjecture), it is natural to expect it does so for general $(r,t)$
as well, since the Macdonald weights only \emph{redistribute} the mass
among cylindric partitions without changing the total count.
\end{remark}


%% ============================================================
\section{Multivariate extensions}
\label{sec:multivariate}
%% ============================================================

The conjecture as stated in Section~\ref{sec:conjecture} uses the
bivariate generating function $\GK_c(z,q)$, which tracks the single
statistic $\max(\mathfrak{c})$ via the variable~$z$.  At Level~2 in the
hierarchy (Section~\ref{sec:hierarchy}), we have the multivariate
generating function with separate variables $z_0, z_1, \ldots, z_{T-1}$
tracking $|\mu^0|, |\mu^1|, \ldots, |\mu^{T-1}|$ individually.
This suggests natural multivariate lifts of the conjecture.

\subsection{Multivariate $E$--$H$ factorisation}

For rank~3 with profile $c = (c_0, c_1, c_2)$, the cylindric partition
$\mathfrak{c} = (\mu^0, \mu^1, \mu^2)$ consists of three partitions.
Define the \emph{multivariate generating function}
\[
  \GK_c(\mathbf{z}, q)
  := \sum_{\mathfrak{c} \in \CPP(c)}
  z_0^{|\mu^0|}\, z_1^{|\mu^1|}\, z_2^{|\mu^2|}\, q^{|\mathfrak{c}|},
\]
where $q^{|\mathfrak{c}|} = q^{|\mu^0| + |\mu^1| + |\mu^2|}$.
(Here $\mathbf{z} = (z_0, z_1, z_2)$ and the $q$-variable is
redundant: $q^{|\mathfrak{c}|} = (z_0 z_1 z_2)^? \cdot q^?$ is
already encoded in the $z$-variables.
More precisely, one can set $q = 1$ and work purely with $z_0, z_1, z_2$,
or set $z_0 = z_1 = z_2 = q$ to recover $\GK_c(q) = \GK_c(1,q)$.)

For the multivariate $E$--$H$ factorisation, define the corrected
generating function
\begin{equation}
  G_c(\mathbf{z})
  := \qpoch{z_0 z_1 z_2}{z_0 z_1 z_2}
  \cdot \GK_c(\mathbf{z}, 1),
  \label{eq:Gc-multi}
\end{equation}
and its multi-index coefficients
\begin{equation}
  g_{\mathbf{n},c}
  := [z_0^{n_0} z_1^{n_1} z_2^{n_2}]\, G_c(\mathbf{z}),
  \qquad
  \mathbf{n} = (n_0, n_1, n_2).
  \label{eq:gn-multi}
\end{equation}

\subsection{The multivariate polynomial $Q_{\mathbf{n},c}$}

\begin{definition}[Multivariate $Q$]
\label{def:Q-multi}
For a multi-index $\mathbf{n} = (n_0, n_1, n_2) \in \N^3$, profile
$c = (c_0, c_1, c_2)$ with $d = c_0 + c_1 + c_2$,
$d \not\equiv 0 \pmod{3}$, define
\begin{equation}
  Q_{\mathbf{n},c}
  := \prod_{k=0}^{2} \qpochn{q}{q}{n_k}
  \cdot g_{\mathbf{n},c}(q),
  \label{eq:Q-multi}
\end{equation}
where $q = z_0 z_1 z_2$ and $g_{\mathbf{n},c}(q)$ is the
appropriately normalised multi-index coefficient.
\end{definition}

\begin{remark}
Setting $n_0 = n_1 = n_2 = n$ and extracting the coefficient of
$z_0^n z_1^n z_2^n$ from the multivariate generating function is
\emph{not} the same as the bivariate $[z^n]$ extraction (which
uses the max statistic).  The relationship between the multivariate
and bivariate conjectures is therefore non-trivial: the bivariate
conjecture concerns a different projection of the same
underlying data.
\end{remark}

\begin{conjecture}[Multivariate positivity]
\label{conj:multi}
Let $c = (c_0, c_1, c_2)$ with $d := c_0 + c_1 + c_2$ and
$d \not\equiv 0 \pmod{3}$.  Then for all
$\mathbf{n} = (n_0, n_1, n_2) \in \N^3$:
\begin{enumerate}
  \item[(i)] $Q_{\mathbf{n},c}$ is a polynomial in~$q$ with
    \textbf{nonneg coefficients}.
  \item[(ii)] Specialising to the diagonal $n_0 = n_1 = n_2 = n$
    and projecting via the max statistic recovers
    Conjecture~\ref{conj:main}.
\end{enumerate}
\end{conjecture}

\subsection{The $(r,t)$-Macdonald multivariate conjecture}

Combining the Macdonald extension (Section~\ref{sec:mac-conjecture})
with the multivariate framework, define the
\emph{Macdonald-weighted multivariate generating function}:
\[
  \GK_c^{(r,t)}(\mathbf{z})
  := \sum_{\mathfrak{c} \in \CPP(c)}
  W_{\mathfrak{c}}(r,t) \cdot
  z_0^{|\mu^0|}\, z_1^{|\mu^1|}\, z_2^{|\mu^2|},
\]
and the corrected multi-index coefficients
\begin{equation}
  g_{\mathbf{n},c}^{(r,t)}
  := [z_0^{n_0} z_1^{n_1} z_2^{n_2}]\!\Bigl(
    \qpoch{z_0 z_1 z_2}{z_0 z_1 z_2} \cdot
    \GK_c^{(r,t)}(\mathbf{z})
  \Bigr).
  \label{eq:gn-multi-rt}
\end{equation}

\begin{definition}[Macdonald multivariate $Q$]
\label{def:Q-multi-rt}
For a multi-index $\mathbf{n} = (n_0, n_1, n_2) \in \N^3$, define
\begin{equation}
  Q_{\mathbf{n},c}(q;\,r,t)
  := \prod_{k=0}^{2} \qpochn{q}{q}{n_k}
  \cdot g_{\mathbf{n},c}^{(r,t)}(q),
  \label{eq:Q-multi-rt}
\end{equation}
where $q = z_0 z_1 z_2$.
\end{definition}

\begin{conjecture}[Macdonald multivariate positivity]
\label{conj:multi-mac}
Let $c = (c_0, c_1, c_2)$ with $d := c_0 + c_1 + c_2$ and
$d \not\equiv 0 \pmod{3}$.  Then for all
$\mathbf{n} = (n_0, n_1, n_2) \in \N^3$:
\begin{enumerate}
  \item[(i)] $Q_{\mathbf{n},c}(q;\,r,t)$ is a polynomial in~$q$
    whose coefficients are rational functions of~$r,t$ that are
    \textbf{nonneg} for $0 < r, t < 1$.
  \item[(ii)] Specialising $r = t$ recovers
    Conjecture~\ref{conj:multi}:
    \[
      Q_{\mathbf{n},c}(q;\,r,r)
      = Q_{\mathbf{n},c}(q).
    \]
  \item[(iii)] Restricting to the bivariate setting (single index
    $n$ via the max statistic) recovers
    Conjecture~\ref{conj:mac}.
\end{enumerate}
\end{conjecture}

\begin{remark}[Chain of implications]
The three conjectures form a chain of increasing generality:
\[
  \text{Conjecture~\ref{conj:multi-mac}}
  \;\;\underset{r=t}{\Longrightarrow}\;\;
  \text{Conjecture~\ref{conj:multi}}
  \;\;\underset{\text{max proj.}}{\Longrightarrow}\;\;
  \text{Conjecture~\ref{conj:main}}
\]
and the bivariate Macdonald version fits as:
\[
  \text{Conjecture~\ref{conj:multi-mac}}
  \;\;\underset{\text{max proj.}}{\Longrightarrow}\;\;
  \text{Conjecture~\ref{conj:mac}}
  \;\;\underset{r=t}{\Longrightarrow}\;\;
  \text{Conjecture~\ref{conj:main}}.
\]
\end{remark}


%% ============================================================
\section{Cylindric Kostka polynomials}
\label{sec:cyl-kostka}
%% ============================================================

The polynomials $Q_{n,c}(q)$ bear a striking structural resemblance to
the classical Kostka--Foulkes polynomials $K_{\lambda,\mu}(t)$.
We argue that they should be regarded as \emph{cylindric} analogues
of Kostka polynomials, and that this perspective clarifies both
the conjecture and the path toward its proof.

\subsection{Recollection: Kostka polynomials}

The \emph{Kostka--Foulkes polynomials} $K_{\lambda,\mu}(t)$
\cite{Macdonald} arise as the transition coefficients between the Schur
and Hall--Littlewood bases of the ring of symmetric functions:
\begin{equation}
  s_\lambda(X)
  = \sum_\mu K_{\lambda,\mu}(t)\, P_\mu(X;t),
  \label{eq:KF-def}
\end{equation}
where $P_\mu(X;t)$ is the Hall--Littlewood polynomial.
Equivalently, inverting the relation,
$P_\mu(X;t) = \sum_\lambda \widetilde{K}_{\lambda,\mu}(t)\, s_\lambda(X)$
expresses the Hall--Littlewood polynomial in the Schur basis.

Their key properties are:
\begin{enumerate}
  \item[(K1)] \textbf{Polynomiality:}
    $K_{\lambda,\mu}(t) \in \Z[t]$.
  \item[(K2)] \textbf{Positivity:}
    $K_{\lambda,\mu}(t)$ has nonneg integer coefficients
    (Lascoux--Sch\"utzenberger \cite{LS}).
  \item[(K3)] \textbf{Combinatorial interpretation:}
    $K_{\lambda,\mu}(t) = \sum_{T} t^{\operatorname{charge}(T)}$,
    summed over semistandard Young tableaux $T$ of shape $\lambda$
    and weight $\mu$, where $\operatorname{charge}$ is Lascoux and
    Sch\"utzenberger's charge statistic.
  \item[(K4)] \textbf{Evaluation:}
    $K_{\lambda,\mu}(1) = K_{\lambda,\mu}$, the Kostka number
    counting semistandard Young tableaux of shape $\lambda$ and
    weight $\mu$.
  \item[(K5)] \textbf{Macdonald lift:}
    There exist two-parameter polynomials $K_{\lambda,\mu}(r,t)$
    (the \emph{Macdonald--Kostka polynomials} or
    $(r,t)$-Kostka polynomials)
    with $K_{\lambda,\mu}(0,t) = K_{\lambda,\mu}(t)$.
  \item[(K6)] \textbf{Representation theory:}
    The Kostka polynomials encode multiplicities in the decomposition
    of representations of $\mathrm{GL}_n$ and appear as
    Kazhdan--Lusztig polynomials for the affine Weyl group of type
    $\mathrm{A}_1^{(1)}$.
\end{enumerate}

\subsection{The cylindric analogue}

The polynomials $Q_{n,c}(q)$ share each of these properties
(proved or conjectured):
\begin{enumerate}
  \item[(C1)] \textbf{Polynomiality:}
    $Q_{n,c}(q) \in \Z[q]$ (proved by Welsh \cite{Welsh}).
  \item[(C2)] \textbf{Positivity:}
    $Q_{n,c}(q)$ has nonneg integer coefficients
    (Conjecture~\ref{conj:main}, open).
  \item[(C3)] \textbf{Combinatorial interpretation:}
    Unknown.  Finding a charge-like statistic on cylindric
    partitions (or cylindric tableaux) is the central open problem.
  \item[(C4)] \textbf{Evaluation:}
    $Q_{n,c}(1) = \bigl(\frac{(d+1)(d+2)}{6} - 1\bigr)^n$
    (proved by Welsh \cite{Welsh}).
  \item[(C5)] \textbf{Macdonald lift:}
    The $(r,t)$-extended polynomial $Q_{n,c}(q;\,r,t)$
    of Definition~\ref{def:Q-rt}, with
    $Q_{n,c}(q;\,r,r) = Q_{n,c}(q)$.
  \item[(C6)] \textbf{Representation theory:}
    The $Q_{n,c}(q)$ appear as numerators in $\mathrm{A}_2^{(1)}$
    Andrews--Gordon characters via equation~\eqref{eq:AG-from-Q}.
\end{enumerate}

\begin{definition}[Cylindric Kostka polynomial]
\label{def:cyl-kostka}
We call $Q_{n,c}(q)$ the \emph{cylindric Kostka polynomial} of
index $n$ and profile $c$.  More generally,
$Q_{\mathbf{n},c}(q;\,r,t)$ (Definition~\ref{def:Q-multi-rt})
is the \emph{$(r,t)$-cylindric Kostka polynomial} of multi-index
$\mathbf{n}$ and profile $c$.
\end{definition}

\subsection{Schur expansion of the Borodin product}

The multivariate Borodin product (Level~2, Section~\ref{sec:hierarchy})
is a formal power series in $z_0, z_1, z_2$:
\[
  \GK_c(\mathbf{z})
  = \sum_{\mathfrak{c} \in \CPP(c)}
  z_0^{|\mu^0|}\, z_1^{|\mu^1|}\, z_2^{|\mu^2|}
  = \prod \text{(Borodin factors)}.
\]
This product can be expanded in terms of Schur functions
$s_\lambda(z_0, z_1, z_2)$ (for partitions $\lambda$ with at most
3~parts):
\begin{equation}
  \GK_c(\mathbf{z})
  = \sum_\lambda a_{\lambda,c}\, s_\lambda(z_0, z_1, z_2),
  \label{eq:schur-expansion}
\end{equation}
where the coefficients $a_{\lambda,c}$ are formal power series.

\begin{remark}
The multivariate generating function $\GK_c(\mathbf{z})$ is
\emph{not} symmetric in $(z_0, z_1, z_2)$ in general (the
profile~$c$ breaks the symmetry unless $c_0 = c_1 = c_2$).
The Schur functions $\{s_\lambda\}_{\ell(\lambda) \leq 3}$
form a basis for the \emph{symmetric} polynomials in three
variables, so the expansion \eqref{eq:schur-expansion} captures
the symmetric part.  The full expansion requires additional
terms from the non-symmetric sector---for instance, Demazure
characters (key polynomials) or the monomial basis.
For a cyclic profile ($c_0 = c_1 = c_2$), the generating
function is symmetric and the Schur expansion suffices.
\end{remark}

The coefficients $a_{\lambda,c}$ in \eqref{eq:schur-expansion}
are not known in closed form.  One might hope that non-negativity
of these Schur expansion coefficients would imply the
positivity conjecture for $Q_{n,c}(q)$.  However, the relationship
between the two is \textbf{not straightforward}, for two reasons:
\begin{enumerate}
  \item \textbf{Max extraction.}
    The bivariate generating function $\GK_c(z,q)$ is obtained
    from the multivariate one by the non-linear operation
    $z_0^{|\mu^0|} z_1^{|\mu^1|} z_2^{|\mu^2|}
    \mapsto z^{\max_k \mu^k_1}\, q^{|\mathfrak{c}|}$.
    This is \emph{not} a linear map on formal power series,
    so Schur-positivity of $\GK_c(\mathbf{z})$ does not directly
    transfer to the bivariate setting.
  \item \textbf{$E$--$H$ factorisation.}
    Even after passing to the bivariate $\GK_c(z,q)$, the
    construction of $Q_{n,c}(q)$ involves multiplication by
    $\qpoch{zq}{q} = \sum_{j \geq 0} (-1)^j q^{j(j+1)/2}
    z^j/(q)_j$, which introduces \textbf{alternating signs}.
    The conjecture asserts that these alternating-sign
    cancellations resolve perfectly into non-negativity after
    clearing denominators---but this cancellation is the
    essential difficulty, and it is not implied by any
    positivity of the Schur expansion coefficients $a_{\lambda,c}$.
\end{enumerate}

In other words, the Schur expansion coefficients of the Borodin
product and the cylindric Kostka polynomials $Q_{n,c}(q)$ are
\textbf{distinct} objects.  The name ``cylindric Kostka polynomial''
for $Q_{n,c}(q)$ is justified not because $Q_{n,c}(q)$ \emph{is}
a Schur expansion coefficient, but because it plays the same
\emph{structural role}: it is the polynomial numerator that appears
when an affine Lie algebra character is decomposed as
$\sum Q_{n,c}(q)/(q)_n$.

\begin{remark}[The $\mathrm{A}_1 \to \mathrm{A}_2$ step]
In the $\mathrm{A}_1$ setting (rank~2, one profile entry), the
Andrews--Gordon identity has the form
$\sum_n q^{f(n)}/(q)_n = \text{product}$, where $f(n)$ is a
quadratic form.  The ``Kostka polynomial'' is the monomial
$q^{f(n)}$---trivially positive.

The $\mathrm{A}_2$ step replaces this monomial by a genuine
polynomial $Q_{n,c}(q)$ with (conjecturally) many terms, all
positive.  This is the essential new phenomenon: at rank~3,
the cylindric Kostka polynomials are \emph{not} monomials
(except for $d = 2$, which reduces to the $\mathrm{A}_1$ case).
\end{remark}

\subsection{The missing charge statistic}

The single most important open problem is finding the combinatorial
counterpart of property~(C3): a \emph{charge-like statistic} on some
set of cylindric objects.

\begin{problem}
\label{prob:charge}
Find a finite set $\mathcal{T}_{n,c}$ with
$|\mathcal{T}_{n,c}| = Q_{n,c}(1)
= \bigl(\frac{(d+1)(d+2)}{6} - 1\bigr)^n$ and a statistic
$\operatorname{ch}\colon \mathcal{T}_{n,c} \to \N$ such that
\begin{equation}
  Q_{n,c}(q) = \sum_{T \in \mathcal{T}_{n,c}} q^{\operatorname{ch}(T)}.
  \label{eq:charge-formula}
\end{equation}
\end{problem}

The natural candidates for $\mathcal{T}_{n,c}$ include:
\begin{itemize}
  \item \textbf{Cylindric tableaux} --- semistandard fillings of a
    cylindric shape, generalising the SSYT that index ordinary Kostka
    polynomials.
  \item \textbf{Paths in the alcove geometry} --- the $\mathrm{A}_2^{(1)}$
    alcove walk model, which provides the representation-theoretic
    setting for the conjecture.
  \item \textbf{Cylindric partitions with bounded max} --- the set
    $\CPP(c)_{\max = n}$ has the right generating function (before
    clearing denominators), and the cyclic descent injection studied
    in \cite{Langer-thesis} provides a partial weight-raising
    bijection on this set.
\end{itemize}

Solving Problem~\ref{prob:charge} would simultaneously prove
Conjecture~\ref{conj:main} and provide a manifestly positive
multisum formula for every $\mathrm{A}_2$ Andrews--Gordon identity.


%% ============================================================
\section{Computational evidence}
\label{sec:evidence}
%% ============================================================

The conjecture has been verified computationally for:
\begin{center}
\begin{tabular}{c|c|l|c}
$d$ & $Q_{n,c}(1)$ & Profiles tested & Max $n$ verified \\ \hline
$1$ & $1^n = 1$ & $(1,0,0)$ & all $n$ (trivial) \\
$2$ & $2^n$ & $(2,0,0)$, $(1,1,0)$ & all $n$ ($Q$ is monomial) \\
$4$ & $4^n$ & $(2,2,0)$, $(2,1,1)$ & $n \leq 3$ \\
$5$ & $6^n$ & $(3,1,1)$, $(3,2,0)$, $(2,2,1)$ & $n \leq 3$ \\
$7$ & $11^n$ & $(3,2,2)$, $(2,3,2)$, $(4,2,1)$ & $n \leq 2$ \\
$8$ & $14^n$ & $(3,3,2)$, $(4,3,1)$, $(4,2,2)$ & $n \leq 1$ \\
\end{tabular}
\end{center}

In every case, all coefficients of $Q_{n,c}(q)$ are nonneg.
Moreover, the corrected coefficients $h_{m,c}(q) = \qpochn{q}{q}{m} \cdot g_{m,c}(q)$
are themselves nonneg polynomials, which is a strictly stronger observation.


\begin{thebibliography}{99}
\bibitem{Andrews}
  G.~E.~Andrews,
  \emph{An analytic generalization of the Rogers--Ramanujan identities for odd moduli},
  Proc.\ Nat.\ Acad.\ Sci.\ USA \textbf{71} (1974), 4082--4085.

\bibitem{ASW}
  G.~E.~Andrews, A.~Schilling, S.~O.~Warnaar,
  \emph{An $\mathrm{A}_2$ Bailey lemma and Rogers--Ramanujan-type identities},
  J.~Amer.\ Math.\ Soc.\ \textbf{12} (1999), 677--702.

\bibitem{Borodin}
  A.~Borodin,
  \emph{Periodic Schur process and cylindric partitions},
  Duke Math.~J.\ \textbf{140} (2007), 391--468.

\bibitem{CDU}
  S.~Corteel, J.~Dousse, A.~K.~Uncu,
  \emph{Cylindric partitions and some new $\mathrm{A}_2$ Rogers--Ramanujan identities},
  Proc.\ Amer.\ Math.\ Soc.\ \textbf{150} (2022), 481--497.
  arXiv:2011.12828.

\bibitem{Corteel}
  S.~Corteel,
  \emph{Rogers--Ramanujan identities and the Robinson--Schensted--Knuth correspondence},
  Proc.\ Amer.\ Math.\ Soc.\ \textbf{145} (2017), 2011--2022.

\bibitem{CW}
  S.~Corteel, T.~Welsh,
  \emph{The $\mathrm{A}_2$ Rogers--Ramanujan identities revisited},
  Ann.\ Combin.\ \textbf{23} (2019), 683--694.

\bibitem{FW}
  O.~Foda, T.~A.~Welsh,
  \emph{Cylindric partitions, $W_r$ characters and the Andrews--Gordon--Bressoud identities},
  J.\ Phys.\ A \textbf{49} (2016), 164004.

\bibitem{GK}
  I.~M.~Gessel, C.~Krattenthaler,
  \emph{Cylindric partitions},
  Trans.\ Amer.\ Math.\ Soc.\ \textbf{349} (1997), 429--479.

\bibitem{Gordon}
  B.~Gordon,
  \emph{A combinatorial generalization of the Rogers--Ramanujan identities},
  Amer.\ J.\ Math.\ \textbf{83} (1961), 393--399.

\bibitem{RR}
  L.~J.~Rogers, S.~Ramanujan,
  \emph{Proof of certain identities in combinatory analysis},
  Cambr.\ Phil.\ Soc.\ Proc.\ \textbf{19} (1919), 211--216.

\bibitem{W}
  S.~O.~Warnaar,
  \emph{The $\mathrm{A}_2$ Andrews--Gordon identities and cylindric partitions},
  arXiv:2111.07550, 2021.

\bibitem{Welsh}
  A.~V.~Welsh,
  cited in \cite{W}, proof of polynomiality and $q=1$ evaluation.

\bibitem{Langer-thesis}
  R.~Langer,
  \emph{Cylindric Partitions and Bijective Proofs},
  PhD thesis, University of Melbourne.

\bibitem{LS}
  A.~Lascoux, M.-P.~Sch\"utzenberger,
  \emph{Sur une conjecture de H.\,O.~Foulkes},
  C.~R.\ Acad.\ Sci.\ Paris \textbf{286A} (1978), 323--324.

\bibitem{Macdonald}
  I.~G.~Macdonald,
  \emph{Symmetric Functions and Hall Polynomials},
  2nd ed., Oxford University Press, 1995.
\end{thebibliography}

\end{document}
